% Generated draft: requires mathematical and source review.
\documentclass[11pt]{article}
\usepackage[margin=25mm]{geometry}
\usepackage{fontspec}
\setmainfont{Latin Modern Roman}
\usepackage{amsmath,amssymb,mathtools,unicode-math}
\setmathfont{Latin Modern Math}
\usepackage{enumitem,needspace,xurl,hyperref}
\setlength{\parindent}{0pt}
\setlength{\parskip}{5pt}
\setlength{\emergencystretch}{3em}
\newcommand{\N}{\mathbb{N}}
\newcommand{\Z}{\mathbb{Z}}
\newcommand{\Q}{\mathbb{Q}}
\newcommand{\R}{\mathbb{R}}
\newcommand{\C}{\mathbb{C}}
\newcommand{\F}{\mathbb{F}}
\newcommand{\A}{\mathbb{A}}
\newcommand{\PP}{\mathbb{P}}
\newcommand{\E}{\mathbb{E}}
\newcommand{\eps}{\varepsilon}
\newcommand{\dd}{\,\mathrm{d}}
\begin{document}
\section*{UnsolvedMath ID: 700837}
\textbf{Source problem code:} catalogue entry
\subsection{Definitions and mathematical statement}
An abelian group \(G\), written additively, is \emph{periodic} if every element has finite order; equivalently,
\[
\forall g\in G\;\exists n\in\Z_{\geq 1}\quad ng=0.
\]
The supplied problem, attributed to S. S. Goncharov, asks:
\begin{quote}
Find a description of autostable periodic abelian groups.
\end{quote}
Thus the requested task is to characterize the periodic abelian groups that are autostable in the sense intended by the source.

\textbf{Unresolved qualification.} The supplied record does not define \emph{autostable}. Consequently, it does not determine the additional property to be characterized or any hypotheses implicit in that terminology. No countability, computability, or presentation assumptions are supplied, and none are imposed here. A fully precise mathematical formulation requires the intended definition of autostability and its domain of applicability.
\subsection{Short English statement}
Describe the abelian groups in which every element has finite order and which satisfy the source's intended notion of autostability. That notion is undefined in the supplied record, so the classification target requires clarification from the source.
\subsection{Sources}
\begin{itemize}
\item Local UnsolvedMath registry, permanent ID \(700837\), title \emph{Kourovka Notebook Problem 12.17}, exact source code \texttt{catalogue entry}. The supplied registry labels the problem \emph{open}; this label is source metadata, not an independently verified assessment.
\item \emph{The Kourovka Notebook (Unsolved Problems in Group Theory)}, original problem 12.17, 12th issue (1992). Proposed by S. S. Goncharov.
\item The supplied background cites the 21st issue edition (2026), \texttt{arXiv:1401.0300v48}. The supplied collection description names E. I. Khukhro and V. D. Mazurov as editors. These bibliographic details are supplied attribution; the cited source has not been independently consulted. No definition of autostability or separate reference defining it is included in the record.
\end{itemize}
\end{document}
